% Options for packages loaded elsewhere
\PassOptionsToPackage{unicode}{hyperref}
\PassOptionsToPackage{hyphens}{url}
\PassOptionsToPackage{dvipsnames,svgnames,x11names}{xcolor}
%
\documentclass[
  11pt,
]{article}
\usepackage{amsmath,amssymb}
\usepackage{iftex}
\ifPDFTeX
  \usepackage[T1]{fontenc}
  \usepackage[utf8]{inputenc}
  \usepackage{textcomp} % provide euro and other symbols
\else % if luatex or xetex
  \usepackage{unicode-math} % this also loads fontspec
  \defaultfontfeatures{Scale=MatchLowercase}
  \defaultfontfeatures[\rmfamily]{Ligatures=TeX,Scale=1}
\fi
\usepackage{lmodern}
\ifPDFTeX\else
  % xetex/luatex font selection
  \setmainfont[ItalicFont=DejaVu Serif,ItalicFeatures={FakeSlant=0.2}]{DejaVu Serif}
  \setmonofont[]{DejaVu Sans Mono}
\fi
% Use upquote if available, for straight quotes in verbatim environments
\IfFileExists{upquote.sty}{\usepackage{upquote}}{}
\IfFileExists{microtype.sty}{% use microtype if available
  \usepackage[]{microtype}
  \UseMicrotypeSet[protrusion]{basicmath} % disable protrusion for tt fonts
}{}
\makeatletter
\@ifundefined{KOMAClassName}{% if non-KOMA class
  \IfFileExists{parskip.sty}{%
    \usepackage{parskip}
  }{% else
    \setlength{\parindent}{0pt}
    \setlength{\parskip}{6pt plus 2pt minus 1pt}}
}{% if KOMA class
  \KOMAoptions{parskip=half}}
\makeatother
\usepackage{xcolor}
\usepackage[margin=0.9in]{geometry}
\usepackage{longtable,booktabs,array}
\usepackage{calc} % for calculating minipage widths
% Correct order of tables after \paragraph or \subparagraph
\usepackage{etoolbox}
\makeatletter
\patchcmd\longtable{\par}{\if@noskipsec\mbox{}\fi\par}{}{}
\makeatother
% Allow footnotes in longtable head/foot
\IfFileExists{footnotehyper.sty}{\usepackage{footnotehyper}}{\usepackage{footnote}}
\makesavenoteenv{longtable}
\setlength{\emergencystretch}{3em} % prevent overfull lines
\providecommand{\tightlist}{%
  \setlength{\itemsep}{0pt}\setlength{\parskip}{0pt}}
\setcounter{secnumdepth}{-\maxdimen} % remove section numbering
\ifLuaTeX
  \usepackage{selnolig}  % disable illegal ligatures
\fi
\IfFileExists{bookmark.sty}{\usepackage{bookmark}}{\usepackage{hyperref}}
\IfFileExists{xurl.sty}{\usepackage{xurl}}{} % add URL line breaks if available
\urlstyle{same}
\hypersetup{
  colorlinks=true,
  linkcolor={Maroon},
  filecolor={Maroon},
  citecolor={Blue},
  urlcolor={Blue},
  pdfcreator={LaTeX via pandoc}}

\author{}
\date{}

\begin{document}

\hypertarget{posterior-uncovered-set-decisions-and-honest-certificates}{%
\section{Posterior uncovered-set decisions and honest
certificates}\label{posterior-uncovered-set-decisions-and-honest-certificates}}

This note supplies proofs for the new \texttt{stepc} implementation. It
is a supporting theory note, not a claim that the results below
establish a novel contribution. The existing paper, code, and experiment
archives are unchanged. The formal frequentist certificate is restricted
to the synthetic fixed-design Bernoulli experiment; real finite-cohort
count envelopes are descriptive diagnostics.

\hypertarget{objects-assumptions-and-the-two-uncertainties}{%
\subsection{1. Objects, assumptions, and the two
uncertainties}\label{objects-assumptions-and-the-two-uncertainties}}

Let \texttt{V=\{0,...,n-1\}}, with \texttt{n\textgreater{}=1}. A
\textbf{strict tournament} has exactly one of \texttt{i→j} and
\texttt{j→i} for each distinct pair. Vertex \texttt{u} covers \texttt{v}
if \texttt{u→v} and every out-neighbor of \texttt{v} is an out-neighbor
of \texttt{u}. The uncovered set \texttt{UC(T)} consists of vertices
with no coverer. It is equivalently the set of vertices that reach every
other vertex in at most two directed edges. Indeed, if \texttt{u} covers
\texttt{v}, neither \texttt{v→u} nor \texttt{v→w→u} is possible.
Conversely, if \texttt{v} cannot reach \texttt{u} in two edges, the
tournament forces \texttt{u→v}, and forces \texttt{u→w} whenever
\texttt{v→w}; thus \texttt{u} covers \texttt{v}.

The stochastic frequentist target is a fixed matrix \texttt{p}
satisfying \texttt{p\_ij+p\_ji=1} and \texttt{p\_ij\ !=\ 1/2} off
diagonal, with

\[
T(p)_{ij}=\mathbf 1\{p_{ij}>1/2\}.
\]

For each observed unordered pair, the number \texttt{m\_ij} of outcomes
is fixed before seeing outcomes, and its observations are i.i.d.
Bernoulli(\texttt{p\_ij}). Independence between \emph{different} pairs
is not needed for the simultaneous confidence proof. A pair omitted by
this fixed design has \texttt{m\_ij=0}. The code records directed counts
\texttt{W\_ij}, so \texttt{m\_ij=W\_ij+W\_ji}. These assumptions exclude
outcome-dependent stopping, adaptive allocation, and without-replacement
sampling from a fixed empirical cohort. No anytime or finite-population
guarantee is claimed.

The separate Bayesian object is a posterior on tournament orientations.
Write

\[
q_{ij}=\Pr(P_{ij}>1/2\mid D),\qquad q_{ji}=1-q_{ij}.
\]

This is an \textbf{orientation probability}, not the posterior mean
\texttt{E{[}P\_ij\textbar{}D{]}}. Under independent pairwise Beta priors
with positive parameters, factorized pairwise Bernoulli likelihoods give
independent Beta posteriors and

\[
q_{ij}=1-I_{1/2}(a_{ij}+W_{ij},b_{ij}+W_{ji}).
\]

The posterior puts no mass at \texttt{P\_ij=1/2}. More generally,
independent edge orientations must be assumed explicitly when using a
product formula. Shared latent parameters or a learned neural model can
create dependent orientations; their marginal \texttt{q} matrix alone
does not identify the joint tournament law. The union-bound certificate
below survives such dependence.

An uncertainty claim conditional on an assumed posterior is not a
frequentist coverage claim about \texttt{p}. A test on human cohort
records also does not establish population coverage. Repeating a count
envelope on a finite set of ranked records without replacement does not
turn these records into i.i.d. Bernoulli draws conditional on the
cohort. A separate finite-population analysis would have to identify its
estimand, design, and a proved bound before being used.

\hypertarget{simultaneous-fixed-design-intervals}{%
\subsection{2. Simultaneous fixed-design
intervals}\label{simultaneous-fixed-design-intervals}}

\textbf{Proposition 1 (standard Hoeffding union bound).} Let
\texttt{M=n(n-1)/2} and \texttt{0\textless{}δ\textless{}1}. For each
observed \texttt{i\textless{}j}, define

\[
r_{ij}=\sqrt{\frac{\log(2M/\delta)}{2m_{ij}}},\qquad
L_{ij}=\max\{0,\widehat p_{ij}-r_{ij}\},\qquad
U_{ij}=\min\{1,\widehat p_{ij}+r_{ij}\}.
\]

Set \texttt{L\_ji=1-U\_ij}, \texttt{U\_ji=1-L\_ij}; give omitted pairs
\texttt{{[}0,1{]}}. Then

\[
\Pr_p\!\left(\bigcap_{i<j}\{p_{ij}\in[L_{ij},U_{ij}]\}\right)
\ge 1-\delta.
\]

For \texttt{n=1} there are no pair events and coverage is trivial; no
logarithm with \texttt{M=0} is evaluated. Diagonal entries are
conventions \texttt{{[}1/2,1/2{]}} and do not participate in the union
bound.

\textbf{Proof.} For an observed pair, Hoeffding's inequality gives
\texttt{Pr(\textbar{}p\_hat-p\textbar{}\textgreater{}r)\ \textless{}=\ 2\ exp(-2m\ r²)=δ/M}.
Clipping to \texttt{{[}0,1{]}} cannot exclude an admissible \texttt{p}
previously inside the interval. Unobserved intervals always contain
\texttt{p}. A union bound over at most \texttt{M} observed unordered
pairs bounds the failure probability by \texttt{δ}. Reverse-direction
intervals are equivalent events, so they do not consume a second
allocation. No cross-pair independence is used. For the singleton there
is only the empty intersection. \(\square\)

The implementation spends \texttt{δ/M} per pair from the predefined
\texttt{M}, even if only some pairs are observed. It does not replace
\texttt{M} by a number selected after observing outcomes. It uses strict
orientations: \texttt{i→j} is \textbf{known} if
\texttt{L\_ij\textgreater{}1/2} (equivalently
\texttt{U\_ji\textless{}1/2}); every other pair remains unknown. On the
simultaneous event, all known edges are edges of the strict true
tournament. Equality at \texttt{1/2} never licenses an edge. Latent ties
are excluded from the target assumption; this note does not silently
orient them.

This is one fixed comparison design and one fixed \texttt{δ}. Running
the rule at several checkpoints does not produce an anytime guarantee. A
finite union over predeclared checkpoints requires an additional error
allocation; adaptive sampling or stopping requires a proved sequential
confidence construction. The code's \texttt{coverage\_guarantee} flag is
an assumption declaration, not an automatic audit of how counts were
generated.

\hypertarget{logical-completion-bounds-and-witnesses}{%
\subsection{3. Logical completion bounds and
witnesses}\label{logical-completion-bounds-and-witnesses}}

Let \texttt{G} be any partial tournament, and \texttt{C(G)} its
unrestricted strict tournament completions. Define exact individual
membership sets

\[
N(G)=\bigcap_{T\in C(G)} UC(T),\qquad
P(G)=\bigcup_{T\in C(G)} UC(T).
\]

These are coordinatewise necessary and possible winners. A subset of
\texttt{P(G)} need not itself be the UC of any completion. The code
enumerates all \texttt{2\^{}r} completions only when the number
\texttt{r} of unknown edges is at most a supplied cap (20 by default);
this exact small-graph oracle is separate from the certificate proof
below.

Define the constructive inner set \texttt{A(G)} by the following test:
for every \texttt{u!=v}, either \texttt{v→u} is known, or there is
\texttt{w} with both \texttt{v→w} and \texttt{w→u} known. Define the
outer set \texttt{B(G)} by removing \texttt{v} when there is a known
\texttt{u→v} such that, for every \texttt{w} distinct from \texttt{u,v},
either \texttt{w→v} is known or \texttt{u→w} is known.

\textbf{Proposition 2 (sound UC certificate bounds).} For every partial
tournament,

\[
A(G)\subseteq N(G)\subseteq UC(T)\subseteq P(G)\subseteq B(G)
\quad\text{for every }T\in C(G).
\]

The implementation only needs these inclusions; it does not use an
unproved claim of completeness of the constructive tests. If
\texttt{A(G)=B(G)}, this common set equals \texttt{UC(T)} for every
completion.

\textbf{Proof.} Every known path is preserved in every completion. Thus
a vertex passing the inner test reaches every opponent in at most two
edges in every completion, and is a UC member by the two-step
characterization. For a removed vertex \texttt{v}, retain its witness
\texttt{u}. Any completion has \texttt{u→v}. If \texttt{v→w} in this
completion, the alternative known edge \texttt{w→v} cannot hold; hence
the test supplies known \texttt{u→w}. Therefore every out-neighbor of
\texttt{v} is an out-neighbor of \texttt{u}, so \texttt{u} covers
\texttt{v} in every completion. This proves the outer exclusion. The
middle inclusions follow from intersection and union. If the two
endpoints coincide, all intermediate sets coincide. \(\square\)

For an inner member, choose one parent for each vertex at distance two
and retain the direct edges from the root to its known out-neighbors.
Every nonroot vertex has exactly one incoming retained edge; all such
edges point from the root or one of its direct children. This gives a
spanning out-tree of depth at most two, with exactly \texttt{n-1} edges.
Its preservation alone proves the root's UC membership. For an excluded
vertex, the code records its fixed covering witness. The scalable
implementation takes \texttt{O(n³)} time; it is not an adaptive
edge-query algorithm.

\textbf{Corollary 2.1 (frequentist UC bounds).} Under Proposition 1's
assumptions and the strict-\texttt{p} assumption, applying Proposition 2
to the known-edge graph gives

\[
\Pr_p\{A(G(D))\subseteq UC(T(p))\subseteq B(G(D))\}\ge1-\delta.
\]

\textbf{Proof.} On the \emph{single whole-graph simultaneous event},
\texttt{T(p)} is a completion of \texttt{G(D)}. Proposition 2 applies to
that completion. The event has probability at least \texttt{1-δ}. Any
certificate selected after viewing the data is covered by the same
event, so no uncorrected selection of an individual pairwise interval is
used. \(\square\)

If the reported inner and outer bounds coincide, the returned UC set is
correct on this event. This is simultaneous set recovery, not a
statement that each reported member individually has posterior
probability \texttt{1-δ}. For a finite-cohort or adaptively stopped
diagnostic, only Proposition 2's conditional structural statement
remains, with no asserted event probability.

\hypertarget{posterior-certificate-probabilities}{%
\subsection{4. Posterior certificate
probabilities}\label{posterior-certificate-probabilities}}

\textbf{Proposition 3 (conditional posterior lower bound).} Fix data
\texttt{D} and a depth-two spanning out-tree \texttt{C\_v} rooted at
\texttt{v}. Its directed edges specify the required orientations. For
any joint posterior law of orientations with marginals \texttt{q\_e},

\[
\pi_v:=\Pr(v\in UC(T)\mid D)
\ge \max\!\left\{0,1-\sum_{e\in C_v}(1-q_e)\right\}.
\]

If the required orientations are jointly independent under the
posterior,

\[
\pi_v\ge\prod_{e\in C_v}q_e
\ge\max\!\left\{0,1-\sum_{e\in C_v}(1-q_e)\right\}.
\]

\textbf{Proof.} Whenever every required orientation holds, the root has
a path of length at most two to every vertex, so it belongs to the UC.
The probability that at least one required edge fails is at most the sum
of its marginal failure probabilities. Complementing and clipping at
zero proves the first inequality. Under joint independence the
probability that all required edges hold is exactly their product. The
same union bound applied to independent edge events gives the final
inequality. \(\square\)

The tree may be selected using \texttt{D} or the fixed posterior
\texttt{q}; this is a conditional posterior assertion about the selected
deterministic tree. It does not establish repeated-data coverage.
Frequentist selection must instead be protected by simultaneous
confidence information. \texttt{posterior\_certificate\_bound} returns
the product only when \texttt{independent\_edges=True}; otherwise it
returns the dependence-robust union bound. Its automatic tree choice
uses \texttt{q\textgreater{}1/2} and does not claim to find the most
probable tree. The supplied tree is validated as a spanning depth-two
out-tree before computing either bound.

\hypertarget{an-exact-posterior-uc-marginal-and-raoblackwellization}{%
\subsection{5. An exact posterior UC marginal and
Rao--Blackwellization}\label{an-exact-posterior-uc-marginal-and-raoblackwellization}}

The following identity is derived here to specify and test the
estimator. Its novelty is unverified; the act of conditioning and the
variance argument are standard and are not presented as new statistical
theory.

\textbf{Proposition 4 (independent-orientation identity).} Under a fixed
posterior with mutually independent unordered edge orientations, let
\texttt{W=N\_T\^{}+(v)} be the random set of out-neighbors of
\texttt{v}. Each \texttt{w!=v} belongs to \texttt{W} independently with
probability \texttt{q\_vw}. Define

\[
H_v(W)=\prod_{u\in V\setminus(W\cup\{v\})}
\left(1-\prod_{w\in W}q_{uw}\right).
\]

Then

\[
\pi_v=E[H_v(W)\mid D]
=\sum_{W\subseteq V\setminus\{v\}}
\left(\prod_{w\in W}q_{vw}\prod_{u\notin W\cup\{v\}}q_{uv}\right)H_v(W).
\]

\textbf{Proof.} Conditional on \texttt{W}, exactly the vertices outside
\texttt{W∪\{v\}} beat \texttt{v}. Such a vertex \texttt{u} covers
\texttt{v} precisely when \texttt{u→w} for every \texttt{w∈W}. Those
edge orientations remain independent with probabilities \texttt{q\_uw},
because conditioning on \texttt{W} fixes only edges incident to
\texttt{v}. Thus its conditional covering probability is
\texttt{∏\_(w∈W)\ q\_uw}. Covering events for distinct outside vertices
involve disjoint unordered edges, one endpoint in \texttt{W} and one
outside; they are conditionally independent. Multiplying the
probabilities of not covering gives
\texttt{H\_v(W)=Pr(v∈UC(T)\textbar{}W,D)}. Total expectation gives the
first equality. Expanding the independent distribution of \texttt{W}
gives the finite sum. \(\square\)

Empty products matter: if \texttt{W} is empty and
\texttt{n\textgreater{}1}, every outside vertex covers \texttt{v}, so
\texttt{H=0}. If \texttt{W} contains all other vertices, the outside
product is empty and \texttt{H=1}. The singleton also gives
\texttt{H=1}. The finite sum is exponential in \texttt{n-1}; it is an
exact small-\texttt{n} oracle, not a polynomial-time algorithm.

For three vertices \texttt{a,b,c}, direct enumeration gives

\[
\pi_a=q_{ab}q_{ac}+q_{ab}(1-q_{ac})q_{bc}
+(1-q_{ab})q_{ac}(1-q_{bc}).
\]

At all \texttt{q=1/2}, this is \texttt{1/2}. A differentiable STE score
with a different value, such as \texttt{5/8} in the existing surrogate,
therefore cannot be identified with this posterior membership
probability without a separate calibration argument. Smoothness alone
gives no such argument.

\textbf{Proposition 5 (scalar Rao--Blackwell variance reduction).} Draw
independent copies of \texttt{W} under the posterior. The sample mean of
\texttt{H\_v(W)} is unbiased for \texttt{π\_v}. For a single copy,

\[
\operatorname{Var}(H_v(W)\mid D)
=\pi_v(1-\pi_v)-E[H_v(W)(1-H_v(W))\mid D]
\le\pi_v(1-\pi_v).
\]

Both sample-mean variances divide by the number of independent draws.

\textbf{Proof.} Let \texttt{Y=1\{v∈UC(T)\}}. Proposition 4 states
\texttt{H\_v(W)=E{[}Y\textbar{}W,D{]}}. Unbiasedness follows by total
expectation. Total variance gives

\[
\operatorname{Var}(Y\mid D)=\operatorname{Var}(E[Y\mid W,D]\mid D)
+E[\operatorname{Var}(Y\mid W,D)\mid D].
\]

Since \texttt{Y} is Bernoulli, \texttt{Var(Y\textbar{}D)=π\_v(1-π\_v)},
while \texttt{Var(Y\textbar{}W,D)=H\_v(W)(1-H\_v(W))}. Rearranging
proves the formula and the nonnegative difference. Independent averages
divide variances by the draw count. \(\square\)

This is a marginal variance statement. Conditioning separately for
different vertices does not automatically reduce cross-vertex covariance
or the variance of a set-level score. Marginal membership probabilities
alone do not supply the cardinality-dependent moments needed for
expected F1.

\hypertarget{bayes-expected-f1-decoding-and-finite-monte-carlo-regret}{%
\subsection{6. Bayes expected-F1 decoding and finite Monte Carlo
regret}\label{bayes-expected-f1-decoding-and-finite-monte-carlo-regret}}

Let \texttt{S=UC(T)} under any fixed posterior joint tournament law,
possibly dependent. \texttt{S} is nonempty: a maximum-outdegree vertex
cannot be covered, since a coverer would beat it and every one of its
out-neighbors, and have strictly larger outdegree. For an action
\texttt{A⊆V}, define

\[
F(A,S)=\frac{2|A\cap S|}{|A|+|S|},\qquad
U(A)=E[F(A,S)\mid D].
\]

The empty action has utility zero. The unrestricted Bayes action
maximizes \texttt{U(A)} over all subsets of \texttt{V}; it need not be a
UC set realizable by a posterior tournament. Requiring that feasibility
is a different optimization problem, not solved by unrestricted GFM.

\textbf{Proposition 6 (established General F-measure Maximizer).} Define

\[
\Delta_{i,k}=E\!\left[\frac{2\mathbf1\{i\in S\}}{k+|S|}\ \middle|\ D\right]
\quad (i\in V,\;k=1,\ldots,n).
\]

For each \texttt{k}, take the \texttt{k} largest \texttt{Δ\_i,k}, then
select the cardinality whose sum is largest. This gives an unrestricted
Bayes expected-F1 action.

\textbf{Proof.} For \texttt{\textbar{}A\textbar{}=k}, linearity of
expectation gives \texttt{U(A)=Σ\_(i∈A)\ Δ\_i,k}. Among sets of size
\texttt{k}, this sum is maximized by taking the largest \texttt{k}
weights, with arbitrary deterministic tie breaking. Searching all
cardinalities optimizes over all nonempty actions. At least one such
action has positive utility because \texttt{S} is nonempty, so the empty
action cannot improve the maximum. This proof uses no independence of
labels or edges. \(\square\)

This is the established GFM rule of Waegeman et al.~(2014), applied to
UC indicator vectors; it is not a new decision theorem. Finite posterior
draws maximize the empirical posterior objective rather than integrating
it exactly.

\textbf{Proposition 7 (standard uniform Monte Carlo regret).} Fix
\texttt{ε\textgreater{}0} and a positive integer \texttt{B}. Conditional
on fixed \texttt{D} and a fixed posterior, let \texttt{S\_1,...,S\_B} be
i.i.d. posterior UC sets and
\texttt{U\_hat(A)=B\^{}\{-1\}Σ\_b\ F(A,S\_b)}. Let \texttt{A\_hat}
maximize \texttt{U\_hat} and \texttt{A*} maximize \texttt{U}. For
\texttt{0\textless{}δ\_MC\textless{}1}, if

\[
B\ge \frac{\log(2\cdot2^n/\delta_{MC})}{2\epsilon^2},
\]

then with probability at least \texttt{1-δ\_MC} over the posterior
simulation,

\[
\sup_A|U_{hat}(A)-U(A)|\le\epsilon,\qquad
U(A^*)-U(A_{hat})\le2\epsilon.
\]

\textbf{Proof.} Each F1 summand is in \texttt{{[}0,1{]}}. Hoeffding's
inequality bounds the two-sided error for a fixed action by
\texttt{2\ exp(-2Bε²)}. A union bound over at most \texttt{2\^{}n}
actions proves uniform error. On that event,
\texttt{U(A*)\ \textless{}=\ U\_hat(A*)+ε\ \textless{}=\ U\_hat(A\_hat)+ε\ \textless{}=\ U(A\_hat)+2ε}.
The selected action may depend on the very same simulation because the
bound is simultaneous over actions. \(\square\)

For \texttt{η\textgreater{}0}, an alternative bounds the \textbf{derived
weights} directly. Every summand defining \texttt{Δ\_hat\_i,k} is in
\texttt{{[}0,1{]}}, so

\[
B\ge\frac{\log(2n^2/\delta_{MC})}{2\eta^2}
\Longrightarrow\max_{i,k}|\Delta_{hat,i,k}-\Delta_{i,k}|\le\eta
\]

with probability at least \texttt{1-δ\_MC}. For
\texttt{\textbar{}A\textbar{}=k}, weight error causes utility error at
most \texttt{kη\textless{}=nη}, yielding regret at most \texttt{2nη} by
the same argument. This is a standard finite-class concentration
consequence, not a novelty claim. The weight bound is not a bound on raw
cardinality probability-table entries. The Monte Carlo failure
probability concerns posterior integration, whereas \texttt{δ} in
Proposition 1 concerns comparison-data sampling. Neither alone proves
the posterior model is correct or a learned surrogate is calibrated.

\hypertarget{a-finite-mixture-likelihood-ratio-training-identity}{%
\subsection{7. A finite-mixture likelihood-ratio training
identity}\label{a-finite-mixture-likelihood-ratio-training-identity}}

This section specifies a separate stochastic training objective. It is a
supporting finite-state likelihood-ratio identity with unverified
novelty; likelihood-ratio differentiation and independent baselines are
standard.

Let \texttt{c∈\{1,...,C\}} index finitely many mixture components and
let \texttt{e=(i,j)} run over unordered pairs with
\texttt{i\textless{}j}. Let \texttt{x\_e(T)=1} indicate the orientation
\texttt{i→j}. With finite differentiable logits \texttt{α\_c(θ)} and
\texttt{z\_c,e(θ)}, define

\[
\rho_c(\theta)=\frac{\exp(\alpha_c(\theta))}
{\sum_{d=1}^{C}\exp(\alpha_d(\theta))},\qquad
q_{c,e}(\theta)=\sigma(z_{c,e}(\theta)).
\]

Conditional on its sampled component, the edge orientations are
independent:

\[
p_\theta(c,T)=\rho_c(\theta)\prod_e
q_{c,e}(\theta)^{x_e(T)}
(1-q_{c,e}(\theta))^{1-x_e(T)}.
\]

All joint state probabilities are positive. The marginal tournament law
\texttt{p\_θ(T)=Σ\_c\ p\_θ(c,T)} usually has dependent edges. A learned
law of this form is only a posterior when a specified inference model
justifies that interpretation; the training identity does not supply
that justification.

Fix a nonempty reference set \texttt{S*}, independent of \texttt{θ}, and
define the reward and training objective

\[
R(T)=F(UC(T),S^*),\qquad
J(\theta)=\sum_{c,T}p_\theta(c,T)R(T).
\]

\textbf{Proposition 8 (finite-state score identity and unbiased
leave-one-out baseline).} At any finite parameter value where the logits
are differentiable,

\[
\nabla_\theta J(\theta)
=E_\theta[R(T)g_\theta(c,T)],
\]

where the joint-state score is

\[
\begin{aligned}
g_\theta(c,T)
&=\nabla_\theta\log p_\theta(c,T)\\
&=\nabla_\theta\alpha_c-
\sum_d\rho_d\nabla_\theta\alpha_d+
\sum_e(x_e(T)-q_{c,e})\nabla_\theta z_{c,e}.
\end{aligned}
\]

For \texttt{s\textgreater{}=2} independent draws \texttt{(c\_b,T\_b)}
from this same current law, put

\[
b_b=\frac{1}{s-1}\sum_{a\ne b}R(T_a),\qquad
\widehat g=\frac{1}{s}\sum_{b=1}^{s}(R(T_b)-b_b)g_\theta(c_b,T_b).
\]

Then \texttt{E\_θ{[}g\_hat{]}=∇\_θ\ J(θ)}. Rewards and baseline values
must be held constant when differentiating the sampled log-probability
surrogate.

\textbf{Proof.} The state space has \texttt{C·2\^{}M} elements. Since
the sum is finite and \texttt{R(T)} does not depend on \texttt{θ},
differentiation commutes with summation:

\[
\nabla_\theta J=
\sum_{c,T}R(T)\nabla_\theta p_\theta(c,T)
=\sum_{c,T}p_\theta(c,T)R(T)\nabla_\theta\log p_\theta(c,T).
\]

The second equality uses positivity. Differentiating softmax gives
\texttt{∇\ log\ ρ\_c=∇α\_c-Σ\_d\ ρ\_d∇α\_d}. For each Bernoulli factor,
differentiating \texttt{x\ log\ σ(z)+(1-x)log(1-σ(z))} gives
\texttt{(x-σ(z))∇z}. Adding these terms proves the displayed score
expression. Normalization also gives

\[
E_\theta[g_\theta(c,T)]
=\sum_{c,T}\nabla_\theta p_\theta(c,T)
=\nabla_\theta 1=0.
\]

For each \texttt{b}, \texttt{b\_b} is a function only of the other
independent draws, so it is independent of \texttt{g\_θ(c\_b,T\_b)}.
Thus \texttt{E{[}b\_b\ g\_b{]}=E{[}b\_b{]}E{[}g\_b{]}=0}. Each
uncentered reward-score term has expectation \texttt{∇J} by the first
part; averaging the centered terms preserves that expectation. Treating
the baseline or reward as a differentiable part of the surrogate would
introduce additional terms not present in this argument. \(\square\)

Sampling and the scored log probability must describe the same law.
Clipping \texttt{q} in one of them while retaining unclipped logits in
the other would break the stated identity. Stable log-sigmoid evaluation
expresses the same finite logit law without an additional mathematical
clipping operation. The theorem does not assert optimizer convergence or
reduction of gradient variance from every baseline choice. The
leave-one-out baseline is undefined at \texttt{s=1}; using a zero
baseline at that draw count retains the uncentered identity.

\textbf{Corollary 8.1 (component-conditional UC integration).} For each
component, use its oriented \texttt{q\_c} matrix in Proposition 4 and
denote the result by \texttt{H\_v,c(W)}. Then under this mixture law,

\[
\Pr_\theta(v\in UC(T))=
\sum_c\rho_c E[H_{v,c}(W)\mid c]
=E[H_{v,c}(W)].
\]

\textbf{Proof.} Conditional on \texttt{c}, the independent-edge proof of
Proposition 4 applies. Averaging its conditional probability by the law
of total probability proves the first equality and sampling
\texttt{(c,W)} gives the second. In fact
\texttt{H\_v,c(W)=E{[}1\{v∈UC(T)\}\textbar{}c,W{]}}, so the
total-variance proof of Proposition 5 applies to this scalar estimator
as well. No marginal edge independence is required after mixing.
\(\square\)

Using mixture marginal edge probabilities
\texttt{q\_bar\_e=Σ\_c\ ρ\_c\ q\_c,e} inside Proposition 4's product
formula is generally incorrect. Likewise, conditional component
independence does not license a certificate probability
\texttt{∏\_e\ q\_bar\_e}. A fixed certificate instead has event
probability \texttt{Σ\_c\ ρ\_c∏\_e\ q\_c,e}, or can use the
dependence-robust union bound from mixture marginals. Joint GFM weights
require the dependence between membership and UC size; averaging scalar
component membership marginals does not provide those weights.

The training objective \texttt{J} scores a randomly sampled
\textbf{whole UC set} against a fixed reference truth. The final GFM
decoder optimizes the expected F1 of an \textbf{action} against a random
modeled UC set. Evaluation against held-out reference truth is a third
quantity. Optimizing \texttt{J} does not, by this theorem, guarantee
improvement of the final GFM action, held-out F1, calibration, or
learning convergence. These objectives must be named separately in
results.

\hypertarget{prior-art-and-research-claims}{%
\subsection{8. Prior art and research
claims}\label{prior-art-and-research-claims}}

The following primary sources delimit what may reasonably be claimed.

\begin{longtable}[]{@{}
  >{\raggedright\arraybackslash}p{(\columnwidth - 4\tabcolsep) * \real{0.3333}}
  >{\raggedright\arraybackslash}p{(\columnwidth - 4\tabcolsep) * \real{0.3333}}
  >{\raggedright\arraybackslash}p{(\columnwidth - 4\tabcolsep) * \real{0.3333}}@{}}
\toprule\noalign{}
\begin{minipage}[b]{\linewidth}\raggedright
Primary source
\end{minipage} & \begin{minipage}[b]{\linewidth}\raggedright
Established component used here
\end{minipage} & \begin{minipage}[b]{\linewidth}\raggedright
Scope distinction
\end{minipage} \\
\midrule\noalign{}
\endhead
\bottomrule\noalign{}
\endlastfoot
\href{https://jmlr.org/papers/volume15/waegeman14a/waegeman14a.pdf}{Waegeman
et al., JMLR 2014, \emph{On the Bayes-Optimality of F-Measure
Maximizers}}, Theorems 8--9 & GFM uses a quadratic matrix of derived
weights under an arbitrary joint label law. & UC posterior labels are an
application, not new expected-F1 decoding. \\
\href{https://webspace.maths.qmul.ac.uk/felix.fischer/publications/abfo_partial.pdf}{Aziz
et al., JAIR 2015, \emph{Possible and Necessary Winners of Partial
Tournaments}}, Section 4.4, Theorems 5--7 & Individual possible and
necessary UC winners are polynomial-time computable; the two-step
characterization is used. & Whether a specified entire set is a possible
UC is NP-complete. Individual bounds do not solve that feasibility
problem. \\
\href{https://arxiv.org/pdf/2509.09312}{Contet, Grandi, and Mengin, AAAI
2026, \emph{Explaining Tournament Solutions with Minimal Supports}},
Corollaries 4--6 & UC minimal supports are depth-two rooted out-trees
with \texttt{n-1} edges; a smallest support can be computed in
\texttt{O(n²)}. & Depth-two certificate structure is prior art. The
linked primary manuscript identifies itself as the extended AAAI 2026
version. \\
\href{https://papers.nips.cc/paper/2016/file/fccb3cdc9acc14a6e70a12f74560c026-Paper.pdf}{Ramamohan,
Rajkumar, and Agarwal, NIPS 2016, \emph{Dueling Bandits: Beyond
Condorcet Winners to General Tournament Solutions}}, Section 4.2 & UCB
methods target UC/TC/Banks winners and give logarithmic
cumulative-regret guarantees. & A fixed-confidence whole-set or
expected-F1-loss stopping guarantee is a different objective, and would
still require comparison with this work. \\
\end{longtable}

\textbf{Support caveat.} With finite counts and positive independent
Beta prior parameters, each \texttt{q\_ij} lies strictly between zero
and one. Consequently every strict tournament has positive posterior
support. For \texttt{n\textgreater{}=2}, each vertex is in some
supported UC (make it a Condorcet winner) and absent from another (make
it a Condorcet loser covered by another vertex). Thus unrestricted
posterior-support possible winners are all vertices and necessary
winners are empty. Informative confidence-completion sets must be
justified separately; they are not the support of this Beta posterior.

\textbf{Research goal A --- unproved, novelty unverified.} Develop an
adaptive noisy comparison algorithm with a fixed-confidence whole-UC or
F1-regret stopping rule whose instance-dependent cost accounts for
discovering and verifying outcome certificates. A displayed tree alone
gives no bound on discovery cost. The fixed-design intervals in this
note do not establish such a theorem.

\textbf{Research goal B --- unproved, novelty unverified.} Obtain a
differentiable approximation to posterior decision risk or the GFM
cardinality-dependent moments with quantitative approximation and
gradient-error control, including dependent learned edge models where
relevant. The exact scalar identity and Rao--Blackwell estimator here do
not prove this research goal, nor a guarantee for the existing STE
training surrogate.

No claimed novelty, superiority, active-learning sample complexity, or
population-coverage result follows from the software tests. They check
implementation identities and every-completion logical soundness,
including all 729 partial tournaments at \texttt{n=4}, 100 generated
partial tournaments at \texttt{n=5} with every completion, and all 1024
complete tournaments at \texttt{n=5}.

\end{document}
